\subsection{开集}

定义 1. 集合 X 上的拓扑结构（或更简单地，拓扑）是这样的结构：它由 X 的子集所成之集 D 给出，而 D 具有下列性质（称为拓扑结构的公理）：

\((\mathbf{O}_{\mathbf{I}})\) \(\mathfrak{D}\) 中集合的任何并都是 \(\mathfrak{D}\) 中的集合。

\((\mathbf{O}_{\mathbf{II}})\) \(\mathfrak{D}\) 中集合的任何有限交都是 \(\mathfrak{D}\) 中的集合。

\(\mathfrak{D}\) 中的集合称为由 \(\mathfrak{D}\) 在 \(\mathbf{X}\) 上定义的拓扑结构的开集。

定义 2. 拓扑空间是赋有拓扑结构的集合。

拓扑空间的元素常称为点。当在集合 X 上定义了一个拓扑之后，这个集合称为拓扑空间 X 的底集。

公理 \((\mathrm{O}_{\mathrm{I}})\) 特别蕴含：D 的空子集族之并，即空集，属于 D。公理 \((\mathrm{O}_{\mathrm{II}})\) 蕴含：D 的空子集族之交，即集合 X，属于 D。

拓扑的例子。对任一集合 X，由 X 与 \(\varnothing\) 组成的子集所成之集满足公理 \((\mathrm{O}_{\mathrm{I}})\) 与 \((\mathrm{O}_{\mathrm{II}})\)，因此在 X 上定义了一个拓扑。由 X 的所有子集组成的集合 \(\mathfrak{P}(\mathbf{X})\) 亦然：它定义的拓扑是 X 上的离散拓扑，而赋有该拓扑的集合 X 称为离散空间。

设 \((\mathbf{U}_i)_{i\in I}\) 是 \(\mathbf{A}\)（拓扑空间 \(\mathbf{X}\) 的一个子集）的覆盖；如果所有的 \(\mathbf{U}_i\) 在 \(\mathbf{X}\) 中都是开集，则称该覆盖为开覆盖。

定义 3. 拓扑空间 X 到拓扑空间 X' 上的同胚是 X 的拓扑结构到 X' 的拓扑结构上的同构；也就是说，按照一般定义，它是 X 到 X' 上的一个双射，且把 X 的开子集所成之集变换为 X' 的开子集所成之集。

称 X 与 \(X'\) 同胚，如果存在 X 到 \(X'\) 上的一个同胚。

例。如果 X 与 \(X'\) 是两个离散空间，则 X 到 \(X'\) 上的任何双射都是同胚。

由同胚的定义立即得到下列判别准则：双射 \(f\)（由拓扑空间 \(\mathbf{X}\) 到拓扑空间 \(\mathbf{X}'\) 上）是同胚的充分必要条件是：在 \(f\) 之下，\(\mathbf{X}\) 中每个开集的像是 \(\mathbf{X}'\) 中的开集，并且在 \(f\) 之下，\(\mathbf{X}'\) 中每个开集的逆像是 \(\mathbf{X}\) 中的开集。

